\section*{अध्याय \ref{ch:b1:cell-unit-of-life} --- कोशिका: जीवन की इकाई}

\begin{solution}{exo:b1:cell-unit-of-life:1}
सारे \omterm{def:b1:organism-environment:organism}{जीव} \omterm{def:b1:cell-unit-of-life:cell}{कोशिकाओं} से बने हैं (सूक्ष्मदर्शी के नीचे देखा गया हर \omterm{def:b1:body-plans-tissues:tissue}{ऊतक},
पादप हो या जांतव, उन्हें दिखाता है); हर \omterm{def:b1:cell-unit-of-life:cell}{कोशिका} किसी \omterm{def:b1:cell-unit-of-life:cell}{कोशिका} से आती है
(भ्रूणों और शैवाल में देखा गया विभाजन; पाश्चर के फ्लास्क स्वतःजनन को
बाहर करते हैं); और वंशागत पदार्थ विभाजन पर आगे जाता है
(समसूत्री विभाजन भर पीछा किए गए गुणसूत्र)।
\end{solution}

\begin{solution}{exo:b1:cell-unit-of-life:2}
\omterm{def:b1:cell-unit-of-life:prokeuk}{केंद्रक} (अनुपस्थित / उपस्थित), झिल्लीबद्ध अंगक (अनुपस्थित /
उपस्थित), गुणसूत्र (एक वर्तुल / कई रैखिक), आकार
(\qtyrange{0.5}{5}{\micro m} / \qtyrange{10}{100}{\micro m}), पेप्टिडोग्लाइकन
की कोशिका-भित्ति (बैक्टीरिया) / सेलुलोस, काइटिन की या कोई नहीं,
कोशिकाकंकाल (अल्पविकसित / विस्तृत)।
\end{solution}

\begin{solution}{exo:b1:cell-unit-of-life:3}
$\qty{1}{mm}/\qty{1}{\micro m} = 1000$ बैक्टीरिया; $1000/20 = 50$ जांतव
\omterm{def:b1:cell-unit-of-life:cell}{कोशिकाएँ}।
\end{solution}

\begin{solution}{exo:b1:cell-unit-of-life:4}
विभेदन तरंगदैर्घ्य और संख्यात्मक द्वारक पर निर्भर करता है, \omterm{def:b1:cell-unit-of-life:resolution}{आवर्धन} पर
नहीं; लगभग $1000\times$ से आगे प्रतिबिंब विवर्तन-सीमा
(\qty{0.2}{\micro m}) के धुँधलेपन को ही बड़ा करता है, बिना कुछ नया
अलगाए।
\end{solution}

\begin{solution}{exo:b1:cell-unit-of-life:5}
$0.61\times 550/0.65 = \qty{516}{nm}$; और नीले में, \qty{422}{nm}। दो
बैक्टीरिया \qty{500}{nm} की दूरी पर: बस किसी तरह, नीले प्रकाश में;
राइबोसोम \qty{30}{nm} की दूरी पर: नहीं, पंद्रह के गुणक से।
\end{solution}

\begin{solution}{exo:b1:cell-unit-of-life:6}
बैंजनी: मोटी पेप्टिडोग्लाइकन भित्ति जो क्रिस्टल वायलेट--आयोडीन
संकुल थामे रखती है, कोई बाह्य झिल्ली नहीं (ग्राम-धनी)। गुलाबी: किसी
बाह्य झिल्ली के नीचे पतली पेप्टिडोग्लाइकन, रंजक धुल गया, प्रतिरंजित (ग्राम-ऋणी)।
\emph{E. coli} एक ग्राम-ऋणी दंड है: यानी गुलाबी दंड।
\end{solution}

\begin{solution}{exo:b1:cell-unit-of-life:7}
सूत्रकणिका: $v = \frac{2}{9}\times(5\times 10^{-7})^2\times 100\times
9.81\times 10^4/10^{-3} = \qty{5.4e-4}{m/s}$; यानी \qty{92}{s} में
\qty{5}{cm}। $1000\,g$ पर \omterm{def:b1:cell-unit-of-life:prokeuk}{केंद्रक}: $v = \frac{2}{9}\times(4\times
10^{-6})^2\times 100\times 9810/10^{-3} = \qty{3.5e-3}{m/s}$; यानी
\qty{14}{s} में \qty{5}{cm}।
\end{solution}

\begin{solution}{exo:b1:cell-unit-of-life:8}
सक्सिनेट डिहाइड्रोजनेज़ सूत्रकणिकीय है: सूत्रकणिकाओं का \qty{80}{\%}
अवसाद 2 में है; अवसाद 1 वाली \qty{15}{\%} वे सूत्रकणिकाएँ हैं जो बिना
टूटी \omterm{def:b1:cell-unit-of-life:cell}{कोशिकाओं} तथा मलबे में फँसी रह गईं या केंद्रकों के साथ नीचे खिंच
गईं; और अधिप्लव वाली \qty{5}{\%} समांगीकरण के दौरान टूटी सूत्रकणिकाओं से
आती हैं। लैक्टेट डिहाइड्रोजनेज़ कोशिकाद्रव्यी और विलेय है: वह अंतिम
अधिप्लव में ही रह जाती है।
\end{solution}

\begin{solution}{exo:b1:cell-unit-of-life:9}
(क) जीवित \omterm{def:b1:cell-unit-of-life:cell}{कोशिकाओं} पर कला-विपर्यास प्रकाश सूक्ष्मदर्शन (यह प्रक्रिया
मिनटों लेती है और जीवित \omterm{def:b1:cell-unit-of-life:cell}{कोशिकाएँ} माँगती है)। (ख) किसी अभिरंजित काट का
प्रकाश सूक्ष्मदर्शन, या पक्के तौर पर TEM (सूत्रकणिकाएँ प्रकाश की सीमा
पर हैं)। (ग) TEM (छिद्र \qty{100}{nm} के हैं)। (घ) SEM (पृष्ठ का उभार)।
\end{solution}

\begin{solution}{exo:b1:cell-unit-of-life:10}
निर्मल फ्लास्क यह बाहर कर देता है कि हवा स्वयं, या उबाला शोरबा स्वयं,
जीवन पैदा करता है: हवा खुली गरदन से स्वच्छंद भीतर आती है। धूल के संपर्क
के बाद धुँधला फ्लास्क दिखाता है कि गरदन ने जिसे रोक रखा था --- हवा के कण
--- वही \omterm{def:b1:organism-environment:organism}{जीवों} का स्रोत है। दोनों मिलकर सिद्ध करते हैं कि सूक्ष्मजीव
धूल से ढोए गए सूक्ष्मजीवों से ही आते हैं। किसी बंद उबाले फ्लास्क पर यह
आपत्ति बची रह जाती कि बंद करने से वह हवा बाहर हो गई जो किसी ``प्राणशक्ति''
को चाहिए थी; हंस-गरदन हवा भीतर आने देती है और केवल धूल बाहर रखती है।
\end{solution}

\begin{solution}{exo:b1:cell-unit-of-life:11}
$10^5 = 2^{n}$ से $n = 16.6$ द्विगुणन मिलते हैं, यानी \qty{415}{min},
लगभग \qty{7}{h}। प्रति लीटर $10^8\times 10^3\,\mathrm{mL}\times\qty{1}{pg} = \qty{0.1}{g}$।
वृद्धि इसलिए रुकती है कि ऑक्सीजन (जिसे माध्यम प्रति लीटर कुछ ही
मिलीग्राम दे सकता है), कोई खनिज, या अम्लों तथा अपशिष्टों का जमाव सीमाकारी
हो जाता है, ग्लूकोज नहीं।
\end{solution}

\begin{solution}{exo:b1:cell-unit-of-life:12}
हर विमा में दस गुना बड़ी कोई \omterm{def:b1:cell-unit-of-life:cell}{कोशिका} प्रति इकाई आयतन झिल्ली-पृष्ठ का
दसवाँ भाग रखती है; और चूँकि झिल्लियाँ ऊर्जा-रूपांतरण तथा परिवहन का तंत्र
ढोती हैं, वह बड़ी \omterm{def:b1:cell-unit-of-life:cell}{कोशिका} भूखी रह जाती। \omterm{def:b1:cell-unit-of-life:prokeuk}{सुकेंद्रकी} झिल्ली को अपने ही भीतर
मोड़कर उसे कई गुना कर लेता है (अंतर्द्रव्यी जालिका, सूत्रकणिका की भीतरी
झिल्ली, थाइलेकॉइड) और इस तरह ऐसे बंद कक्ष भी बना लेता है जिनमें असंगत
प्रक्रियाएँ --- अम्लीय पाचन, ऑक्सीकारी रसायन, DNA का भंडारण --- अपनी-अपनी
दशाओं पर अलग-अलग चलती हैं। यह चित्र जहाँ तक झिल्लियों की बात है वहाँ तक
सही है, और अधूरा है: \omterm{def:b1:cell-unit-of-life:prokeuk}{सुकेंद्रकी} ने एक कोशिकाकंकाल और एक \omterm{def:b1:cell-unit-of-life:prokeuk}{केंद्रक} भी पाया,
और उसकी सूत्रकणिकाएँ निगले गए बैक्टीरिया के वंशज हैं, मोड़ नहीं।
\end{solution}

\begin{solution}{pb:b1:cell-unit-of-life:1}
\textbf{1.} अवसाद 1: \omterm{def:b1:cell-unit-of-life:prokeuk}{केंद्रक} (और बिना टूटी \omterm{def:b1:cell-unit-of-life:cell}{कोशिकाएँ}, मलबा)। अवसाद 2:
सूत्रकणिकाएँ (लयनकायों तथा परऑक्सीकायों के साथ)। अवसाद 3: सूक्ष्मकाय
(अंतर्द्रव्यी जालिका के टुकड़े) और राइबोसोम।
\textbf{2.} ठंड एंजाइमों को धीमा करती है (उन प्रोटीन-अपघटनियों तथा
फ़ॉस्फ़ोलाइपेज़ों को भी जो टूटने से छूटती हैं) और अंगकों को बचाती है;
और समपरासरी इसलिए कि अंगक न फूलकर फूटें और न सिकुड़ें।
\textbf{3.} सूत्रकणिकाएँ (भीतरी झिल्ली); \omterm{def:b1:cell-unit-of-life:cell}{कोशिकाद्रव्य}।
\textbf{4.} $500 + 200 + 300 + 950 = \qty{1950}{mg}$, \qty{97.5}{\%}:
नलियों की दीवारों पर और स्थानांतरणों में हुई हानि।
\textbf{5.} SDH: $15 + 60 + 5 + 4 = 84$, \qty{84}{\%} (कुछ सक्रियता
सूत्रकणिकाओं की क्षति से खोई)। LDH: $20 + 12 + 8 + 350 = 390$,
\qty{97.5}{\%}।
\textbf{6.} $2000/10 = \qty{200}{mg/g}$: यानी ताज़े द्रव्यमान का \qty{20}{\%} (बाकी अधिकांश जल है)।
\textbf{7.} $v = \frac{2}{9}\times(5\times 10^{-7})^2\times 100\times
\num{9.81e4}/10^{-3} = \qty{5.4e-4}{m/s}$; \qty{92}{s} में \qty{5}{cm}:
\qty{20}{min} बहुत है।
\textbf{8.} $1000\,g$ पर, $v = \qty{5.4e-5}{m/s}$; \qty{600}{s} में,
\qty{3.3}{cm}। \qty{5}{cm} भर एकसमान बँटी हुईं, तली से \qty{3.3}{cm} के
भीतर की सूत्रकणिकाएँ उस तक पहुँच जाती हैं: यानी \qty{65}{\%}, यदि और कुछ
बीच में न आता।
\textbf{9.} सचमुच केवल \qty{15}{\%} अवसादित होती हैं: कोई असली अवसाद उस
सब को नहीं फँसाता जो उस तक पहुँचता है (जमे \omterm{def:b1:cell-unit-of-life:prokeuk}{केंद्रक} धीरे से फिर निलंबित
किए जाते हैं और ढीली ऊपरी परत अधिप्लव में रह जाती है), और सूत्र अगोलीय,
जलयोजित अंगकों की चाल बढ़ा-चढ़ाकर आँकता है।
बिना टूटी \omterm{def:b1:cell-unit-of-life:cell}{कोशिकाओं} के भीतर और केंद्रकों से चिपकी सूत्रकणिकाएँ भी अवसाद 1
में चली जाती हैं।
\textbf{10.} $v = \frac{2}{9}\times(4\times 10^{-6})^2\times 100\times
9810/10^{-3} = \qty{3.5e-3}{m/s}$; \qty{14}{s}।
\textbf{11.} $v = \frac{2}{9}\times(1.2\times 10^{-8})^2\times 600\times
\num{9.81e5}/10^{-3} = \qty{1.9e-5}{m/s}$; \qty{2650}{s} में \qty{5}{cm},
यानी \qty{44}{min}। $\num{10000}\,g$ पर चाल \qty{1.9e-6}{m/s} है:
\qty{20}{min} में \qty{2.3}{mm}, इसलिए राइबोसोम अधिप्लव 2 में रह जाते हैं।
\textbf{12.} चाल $r^2$ की तरह चलती है: नीचे बल पर केवल सबसे बड़े कण ही
ध्यान देने लायक चलते हैं, इसलिए हर चरण एक वर्ग अवसादित करता है और छोटों
को छोड़ देता है। बहुत छोटा रखिए तो वह वर्ग अधूरा अवसादित होता है; बहुत
लंबा तो अगला वर्ग अवसाद को दूषित करने लगता है।
\textbf{13.} \omterm{def:b1:cell-unit-of-life:fractionation}{समांगीकृत} $100/2000 = \qty{0.05}{U/mg}$; अवसाद 1
$15/500 = 0.03$; अवसाद 2 $60/200 = 0.30$; अवसाद 3 $5/300 = 0.017$;
अधिप्लव $4/950 = 0.004$।
\textbf{14.} शुद्धिकरण $0.30/0.05 = 6$; उपज $60/100 =
\qty{60}{\%}$।
\textbf{15.} $0.017/0.05 = 0.33$: अवसाद 3 सूत्रकणिकाओं में दरिद्र है; उसमें जो थोड़ी सक्रियता है वह समांगीकरण से बनी सूत्रकणिका-झिल्ली के टुकड़ों की है, जो सूक्ष्मकायों के साथ अवसादित होते हैं।
\textbf{16.} LDH: \omterm{def:b1:cell-unit-of-life:fractionation}{समांगीकृत} $400/2000 = \qty{0.2}{U/mg}$; अधिप्लव
$350/950 = \qty{0.37}{U/mg}$; शुद्धिकरण $1.8$; उपज \qty{87.5}{\%}।
\textbf{17.} \omterm{def:b1:cell-unit-of-life:cell}{कोशिकाद्रव्य} \omterm{def:b1:cell-unit-of-life:cell}{कोशिका} के प्रोटीन का आधा है, इसलिए कोई शुद्ध
\omterm{def:b1:cell-unit-of-life:cell}{कोशिकाद्रव्य} भी विशिष्ट सक्रियता को दुगना ही कर सकता है; सूत्रकणिकाएँ
छठा भाग या उससे कम हैं, इसलिए उन्हें अलग करने से उनके एंजाइम छह गुना या
अधिक सांद्रित होते हैं। विलेय एंजाइम की उपज ऊँची है क्योंकि वह क्षतिग्रस्त
अंगकों के साथ नहीं खोता।
\textbf{18.} नहीं: अवसाद 2 \qty{200}{mg} प्रोटीन रखता है और अपने कणों के
बीच अधिप्लव फँसाता है; \qty{12}{U} कुल का \qty{3}{\%} है, जो फँसे
\omterm{def:b1:cell-unit-of-life:cell}{कोशिकाद्रव्य} के अनुकूल है। परीक्षण: अवसाद धोइए
(फिर निलंबित कीजिए और फिर घुमाइए): कोई फँसा एंजाइम निकल जाता है, कोई
बँधा हुआ टिका रहता है।
\textbf{19.} यदि अवसाद 2 \qty{60}{\%} सूत्रकणिकाएँ हो, तो शुद्ध
सूत्रकणिकाओं का $0.30/0.6 = \qty{0.5}{U/mg}$ होता है; तब यकृत की
\qty{100}{U} \qty{200}{mg} सूत्रकणिकीय प्रोटीन के अनुरूप हैं, यानी
\qty{2000}{mg} का \qty{10}{\%}।
\textbf{20.} स्थिरित, अंतःस्थापित, पतली काट में काटे, भारी-धातु से
अभिरंजित अवसाद का TEM: दोहरी झिल्ली और क्रिस्टी वाले सेम-जैसे परिच्छेद,
जिनके बीच कुछ लयनकाय (सघन पिंड) और परऑक्सीकाय होते हैं।
\textbf{21.} फिर निलंबित किए अवसाद की एक बूँद कला-विपर्यास सूक्ष्मदर्शी के नीचे देखिए और मुक्त केंद्रकों के बीच अक्षत \omterm{def:b1:cell-unit-of-life:cell}{कोशिकाएँ} गिनिए; यदि बहुत हों, तो अधिक अच्छी तरह \omterm{def:b1:cell-unit-of-life:fractionation}{समांगीकृत} कीजिए (अधिक स्ट्रोक, अधिक कसा मूसल) और फिर घुमाइए, क्योंकि फँसे अंगक हर उपज गिरा देते हैं।
\textbf{22.} विभेदी अपकेंद्रण आकार से अलग करता है (चाल
$\propto r^2$) और तीनों अंगकों के आकार मिलते-जुलते हैं; किसी
घनत्व-प्रवणता में कोई कण अपने ही घनत्व पर रुक जाता है, चाहे उसका आकार
कुछ भी हो, और तीनों के घनत्व भिन्न हैं।
\textbf{23.} अक्षत लयनकाय एंजाइम को ऐसी झिल्ली के भीतर रखते हैं जिससे
क्रियाधार पार नहीं होता (प्रच्छन्नता); अपमार्जक झिल्ली तोड़ देता है और
उसे उजागर कर देता है। इसलिए अवसाद के लयनकाय अधिकतर अक्षत हैं।
\textbf{24.} जल अल्पपरासरी है: अंगक फूलकर फट जाते हैं।
सूत्रकणिकाएँ और लयनकाय फट जाते हैं, इसलिए अवसाद 2 सिकुड़ जाता है;
सक्सिनेट डिहाइड्रोजनेज़ (झिल्ली-बद्ध) अंशतः झिल्ली के टुकड़ों के रूप में
अवसाद 3 में चली जाती है; अम्ल फ़ॉस्फ़ेटेज़ (लयनकाय के भीतर विलेय) अंतिम
अधिप्लव में छूट जाती है।
\textbf{25.} शुद्धिकरण गुणक $6$, उपज \qty{60}{\%}। यह गुणक
संदूषण से सीमित है (लयनकाय, परऑक्सीकाय, फँसा \omterm{def:b1:cell-unit-of-life:cell}{कोशिकाद्रव्य}:
अवसाद 2 \qty{60}{\%} सूत्रकणिकाएँ है); और उपज उन सूत्रकणिकाओं से जो
अवसाद 1 को खो गईं (फँसकर) और समांगीकरण के दौरान टूट गईं।
\end{solution}
